% \noindent \textbf{\Large Argonne National Laboratory}
\phantomsection
\markright{Argonne National Laboratory}

% \addcontentsline{toc}{subsection}{\numberline {} Mathematics and Computer
% Science Division/ALCF}

The lead PI will take responsibility for the collection, management, and sharing of
the research data.  The co-PIs and project members will conform to best practices and
standards to facilitate the reproduction and verification of our results.

\subsection*{Data Types and Sources}

This project will produce two general classes of data: new and expanded software 
and simulation results as a function of problem specification, input 
parameters, hardware specifications, software optimizations, and so on. The 
software will be accompanied by technical documentations. Any simulation 
results this project publishes will be documented with the hardware and software 
information, to facilitate reproduction/validation where possible.

%The project itself will not produce any such sensitive data.
%All project data will be documented on a project wiki: a website
%that all project members may edit.
%
%
%This project will result in new and expanded software
%libraries, all of which will be open source and freely available on the web.
%
%Evaluation of the proposed framework will result in
%experimental data that will be published in either a technical report
%and included in Argonne's publication
%systems, also freely accessible from the web.

\subsection*{Content and Format}

The publications will be stored in Latex, PDF, or PowerPoint (for posters) file
format. The source code of the software and the analysis scripts will be in C,
C++, Python, and R files. Data from simulations will be
stored using CSV, JSON, or flat file formats. Related files in different formats
will be linked by file naming conventions. For each simulation, one
metadata file will be generated and will be stored in the corresponding data set
directory. Dublin Core metadata standard will be applied during the creation of
the metadata.

\subsection*{Sharing and Preservation}
%a.	Means for sharing and rationale behind any restrictions on who may access the data
The project itself will not produce any sensitive data.
The project web page will be hosted at the Argonne National Laboratory, where the
URLs for the publication, software, and simulation data will be indexed with
appropriate headings/titles.
All research data from publications resulting from this project will be made available in an open, machine-readable, and digitally accessible format to the public at
the time of publication, either as electronic supplements, or on the Argonne-hosted project page that will be linked to the publication.
%at the  b.	Timeline for sharing and preservation
All data produced by this project will be retained for the duration of the
project and a minimum of 12 months following the project's completion.
%c.	Special requirements for data sharing
The data sharing does not require special infrastructure.
%d.	Statement of resources and capabilities requested in the research proposal that are needed to meet sharing and preservation
The preprints of publications from Argonne and IIT team members will be made available at
\url{http://www.mcs.anl.gov/publications}. The project webpage, wiki page,
software, and experimental data will be stored
at \url{https://xgitlab.cels.anl.gov/}, %\url{https://repo.anl-external.org/},
hosted by the Argonne National Laboratory. In addition to these repos, we will
use Box, Argonne's secure cloud-based file storage system that provides ways to
organize, edit, share, and collaborate on data. Each user has a quota of 1.5 TB
of data. Aside from public downloads, Box can also be used for private sharing
of collaborative data. The organization hosting the services will utilize best
practices for service providers, including regular backups, disaster recovery
plans, security plans, and general availability guarantees.
%e.	Cost/benefit considerations detailing whether or where the data will be preserved after the direct project funding ends
On completion of significant products (white papers, benchmarks, the project
itself), data will be made public and copies of the data will be copied to DVD
or tapes for long-term retention.
%f.	Conditions under which the research data will be transferred to a third party (institutional or community repository)
%g.	Future decision points regarding reevaluation of costs and benefits of data sharing and preservation

\subsection*{Protection}
This project will not involve any personally identifiable information. However,
a record of the complete state of a leadership-class machines (job placement,
other user's concurrent jobs, power state, etc.) is not feasible.
Any intermediate simulation data or source code implementations produced prior
to the published results will not be shared or preserved.

%\section{Rationale}

\subsection*{Software and Codes}

We will follow Argonne's procedure for releasing software as open source.
The developed software and codes will be stored internally at
\url{https://xgitlab.cels.anl.gov/}, and will be accessible to 
all co-PIs, postdocs, graduate students, and summer interns.
Before making software available to the public, we will first register it with
ANL and make sure that it has an appropriate license and copyright notice.
